\documentclass{article}
\usepackage{amsmath}
\usepackage{amsfonts}
\usepackage{geometry}
\geometry{a4paper, margin=1cm}

\title{Hybrid Algorithm: Compact Description}
\author{}
\date{}

\begin{document}

\maketitle

\section*{Overview}
Universal goal-driven problem-solving algorithm combining:

    
- Resonance analysis
    
- Dynamic constraints
    
- Reinforcement learning (RL + GAN)
    
- Attention mechanisms


\section*{Core Equation}
\[
P_{\text{total}} = 1 - \prod_{i=1}^n (1 - P_i)
\]
Measures total success probability across independent subgoals.

\section*{Key Components}
\subsection*{1. Resonance Analysis}
Finds critical points where small changes yield large effects:
\[
\omega_{\text{res}} = \frac{1}{D} \cdot \sum_{k=1}^N \frac{q_k}{m_k}
\]

\subsection*{2. Dynamic Constraint Manipulation}
Uses tokens like \texttt{[VAR\_c]}, \texttt{[CONST\_c]} to temporarily change constants.

\subsection*{3. Attention Mechanism}
\[
\alpha_i = \frac{e^{\omega_{\text{res},i}}}{\sum_j e^{\omega_{\text{res},j}}}
\]
Focuses attention on high-impact parameters.

\subsection*{4. RL + GAN}

    
- \textbf{Generator}: Proposes new hypotheses
    
- \textbf{Discriminator}: Evaluates feasibility
    
- \textbf{RL}: Reinforces successful strategies


\section*{Complexity}
\begin{tabular}{|l|l|}
\hline
Algorithm & Complexity \\
\hline
Basic & $ O(2^{m+n}) $ \\
Hybrid & $ O(n^2) $ \\
\hline
\end{tabular}

\section*{Advantages}

    
- Handles large search spaces
    
- Learns from data
    
- Generates novel hypotheses
    
- Works with sequences and context


\section*{Use Cases}

    
- Drug discovery
    
- Autonomous driving
    
- Space navigation
    
- Nanotech design


\section*{Limitations}

    
- Requires training data
    
- Needs GPU for training
    
- May generate "impossible" scenarios


\section*{When to Use Basic Algorithm}

    
- Small number of subgoals (3–5)
    
- No GPU/data
    
- Deterministic output needed


\end{document}